\documentclass[12pt, a4paper]{article}
 
% ─── Packages ───────────────────────────────────────────────────────────────
\usepackage[a4paper, top=1in, bottom=1in, left=1.2in, right=1.2in]{geometry}
\usepackage{graphicx}
\usepackage{xcolor}
\usepackage{fancyhdr}
\usepackage{titlesec}
\usepackage{hyperref}
\usepackage{booktabs}
\usepackage{array}
\usepackage{listings}
\usepackage{parskip}
\usepackage{setspace}
\usepackage{enumitem}
\usepackage{float}
\usepackage{mdframed}
\usepackage{tikz}
\usetikzlibrary{shapes.geometric, arrows.meta, positioning, fit}
 
% ─── Colours ────────────────────────────────────────────────────────────────
\definecolor{PrimaryBlue}{HTML}{1A3A5C}
\definecolor{AccentBlue}{HTML}{2E75B6}
\definecolor{LightGray}{HTML}{F5F5F5}
\definecolor{MidGray}{HTML}{AAAAAA}
\definecolor{CodeBg}{HTML}{F0F4F8}
\definecolor{CodeKw}{HTML}{2E75B6}
\definecolor{CodeStr}{HTML}{C0392B}
\definecolor{CodeCmt}{HTML}{27AE60}
\definecolor{PatternGreen}{HTML}{1E6B45}
\definecolor{PatternGreenLight}{HTML}{E8F5EE}
 
% ─── Hyperref ───────────────────────────────────────────────────────────────
\hypersetup{colorlinks=true, linkcolor=AccentBlue, urlcolor=AccentBlue,
  pdftitle={PDF Numerical Data Extractor — OOP Group Project}}
 
% ─── Section Formatting ─────────────────────────────────────────────────────
\titleformat{\section}{\large\bfseries\color{PrimaryBlue}}{\thesection}{1em}{}
  [{\color{MidGray}\titlerule[0.4pt]}]
\titleformat{\subsection}{\normalsize\bfseries\color{AccentBlue}}{\thesubsection}{1em}{}
 
% ─── Code Style ─────────────────────────────────────────────────────────────
\lstdefinestyle{py}{
  backgroundcolor=\color{CodeBg},
  commentstyle=\color{CodeCmt}\itshape,
  keywordstyle=\color{CodeKw}\bfseries,
  stringstyle=\color{CodeStr},
  basicstyle=\ttfamily\footnotesize,
  breaklines=true,
  numbers=left,
  numberstyle=\tiny\color{MidGray},
  numbersep=8pt,
  frame=single,
  rulecolor=\color{MidGray},
  tabsize=4,
  showstringspaces=false,
  language=Python,
  morekeywords={self,True,False,None,abstractmethod,ABC},
}
\lstset{style=py}
 
% ─── Info / Pattern Boxes ───────────────────────────────────────────────────
\mdfdefinestyle{infobox}{backgroundcolor=LightGray, linecolor=AccentBlue,
  linewidth=2pt, leftline=true, rightline=false, topline=false, bottomline=false,
  innerleftmargin=10pt, innerrightmargin=10pt, innertopmargin=6pt, innerbottommargin=6pt}
 
\mdfdefinestyle{patternbox}{backgroundcolor=PatternGreenLight,
  linecolor=PatternGreen, linewidth=2pt, leftline=true, rightline=false,
  topline=false, bottomline=false,
  innerleftmargin=10pt, innerrightmargin=10pt, innertopmargin=6pt, innerbottommargin=6pt}
 
% ─── Header / Footer ────────────────────────────────────────────────────────
\pagestyle{fancy}\fancyhf{}
\fancyhead[L]{\small\color{AccentBlue}\textbf{PDF Numerical Data Extractor}}
\fancyhead[R]{\small\color{MidGray}OOP Group Project}
\fancyfoot[C]{\small\color{MidGray}\thepage}
\renewcommand{\headrulewidth}{0.5pt}
\renewcommand{\headrule}{\hbox to\headwidth{%
  \color{AccentBlue}\leaders\hrule height\headrulewidth\hfill}}
 
\setstretch{1.3}
\setlength{\parskip}{5pt}
 
% ════════════════════════════════════════════════════════════════════════════
\begin{document}
% ════════════════════════════════════════════════════════════════════════════
 
% ─── Title Page ─────────────────────────────────────────────────────────────
\begin{titlepage}
\centering
\vspace*{1cm}
{\color{PrimaryBlue}\rule{\linewidth}{2pt}}\\[0.5cm]
 
\includegraphics[width=3cm]{jadavpurlogo.jpg}\\[0.5cm]
 
{\color{MidGray}\small Department of Computer Science \& Engineering}\\[0.1cm]
{\color{MidGray}\small Academic Year 2024--2025}\\[0.6cm]
{\color{PrimaryBlue}\rule{\linewidth}{0.5pt}}\\[0.5cm]
 
{\color{PrimaryBlue}\Huge\bfseries PDF Numerical\\[0.2cm]Data Extractor}\\[0.3cm]
{\color{AccentBlue}\large\itshape An Object-Oriented Pipeline using Factory,\\
Decorator, and Polymorphism Design Patterns}\\[0.7cm]
 
{\color{PrimaryBlue}\rule{\linewidth}{0.5pt}}\\[0.6cm]
 
\begin{minipage}{0.45\textwidth}
\centering{\large\bfseries\color{PrimaryBlue}Group Members}\\[0.4cm]
\begin{tabular}{ll}
\toprule
\textbf{Name} & \textbf{Roll No.}\\
\midrule
Member 1     & XXXXXX01 \\
Member 2     & XXXXXX02 \\
Member 3     & XXXXXX03 \\
Member 4     & XXXXXX04 \\
Member 5     & XXXXXX05 \\
\bottomrule
\end{tabular}
\end{minipage}
\hfill
\begin{minipage}{0.45\textwidth}
\centering{\large\bfseries\color{PrimaryBlue}Project Details}\\[0.4cm]
\begin{tabular}{ll}
\toprule
\textbf{Field} & \textbf{Value}\\
\midrule
Course     & CS2XX: OOP Lab \\
Semester   & 4th / Even \\
Group No.  & 5 \\
Supervisor & Prof.\ [Name] \\
\bottomrule
\end{tabular}
\end{minipage}
 
\vfill
{\color{PrimaryBlue}\rule{\linewidth}{2pt}}\\[0.3cm]
{\color{MidGray}\small Submitted in partial fulfilment of the requirements
for the Object-Oriented Programming Laboratory}
\end{titlepage}
 
% ════════════════════════════════════════════════════════════════════════════
\section*{Abstract}
% ════════════════════════════════════════════════════════════════════════════
 
This report describes the design and implementation of a Python-based pipeline
that solves Question~5 of the OOP group project: given a scanned or printed
PDF, the system extracts all numerical information, converts numbers written in
words into digit form, and classifies each by type across an extended
24-label taxonomy (Date, Money, Percent, Phone Number, License Plate, Time,
Age, Temperature, Distance, Weight, Height, Speed, Area, Volume,
Percentage Change, Ratio, Range, Ordinal, PIN Code, Account Number,
Invoice ID, Tax ID, Quantity, Cardinal). The output is a structured JSON file per document
containing the \emph{original text}, \emph{normalised value}, and
\emph{category} of every detected number.
 
The implementation spans a modular \texttt{components/} package and a
Streamlit front-end. Each module demonstrates a distinct OOP design pattern:
\textbf{Factory Method} (extractor), \textbf{Decorator} (normalisation), and
\textbf{Polymorphism with Liskov Substitution} (parser/classifier), with
output formatting delegated to a dedicated serializer.
 
% ════════════════════════════════════════════════════════════════════════════
\section{Introduction}
% ════════════════════════════════════════════════════════════════════════════
 
\begin{mdframed}[style=infobox]
\textbf{Problem Statement (Q5):} Design a Python function or model that
extracts all numerical information from a scanned or printed PDF, converts
numbers written in words into digits, and classifies each by type.\\[4pt]
\textbf{Input:} A PDF --- scanned or printed.\quad
\textbf{Output:} A JSON file with \emph{original text}, \emph{normalised
value}, and \emph{category} for every number found.
\end{mdframed}
 
Numerical data embedded in PDFs --- dates on invoices, phone numbers on
contact sheets, amounts on government forms --- is inaccessible to automated
systems unless explicitly extracted. This project automates that process
end-to-end, handling both digital and scanned documents without requiring
any user configuration, and outputting results in the JSON schema required
by the assignment.


% ════════════════════════════════════════════════════════════════════════════
\section{Project File Structure}
% ════════════════════════════════════════════════════════════════════════════
 
\begin{lstlisting}[language=bash, numbers=none, caption={Project directory layout}]
oops_project_4thsem/
  app.py                     # Streamlit UI layer (client)
  requirements.txt           # Reproducible pinned dependencies
  components/
    __init__.py
    extractor.py             # Document ingestion engine  [Factory Method]
    normalization.py         # Text pre-processing engine [Decorator Pattern]
    parser.py                # Classification engines     [Polymorphism / LSP]
    serialize.py             # JSON serializer
    main.py                  # ProcessingPipeline orchestration service
\end{lstlisting}
 
Each file owns a focused responsibility. \texttt{app.py} handles UI concerns,
while the orchestration logic is encapsulated in
\texttt{components/main.py} as \texttt{ProcessingPipeline}.

% ════════════════════════════════════════════════════════════════════════════
\section{Workflow Diagram}
% ════════════════════════════════════════════════════════════════════════════
\begin{figure}[H]
  \centering
  \includegraphics[width=0.8\linewidth]{workflow_diagram.png}
  \caption{End-to-end pipeline architecture and OOP design patterns.}
  \label{fig:workflow}
\end{figure}
% ════════════════════════════════════════════════════════════════════════════
\section{Module Design}
% ════════════════════════════════════════════════════════════════════════════
 
% ── 3.1 ─────────────────────────────────────────────────────────────────────
\subsection{\texttt{extractor.py} --- Ingestion Layer (Factory Method Pattern)}
 
\textbf{What it does.} Determines whether the uploaded PDF has a native text
layer (digital) or consists of scanned images (requires OCR), then extracts
the raw text accordingly using \texttt{PyMuPDF} or \texttt{Tesseract}.
 
\textbf{Design pattern: Factory Method.} An interface defines the contract;
two concrete classes implement it; a factory decides at runtime which to
instantiate.
 
\begin{mdframed}[style=patternbox]
\textbf{OOP Principle --- Open/Closed (OCP).} The factory is \emph{open for
extension} (add a \texttt{WordDocProcessor} by writing one new class) but
\emph{closed for modification} (the existing PDF logic is never touched).
\end{mdframed}
 
\begin{lstlisting}[caption={extractor.py --- interface and factory skeleton}]
from abc import ABC, abstractmethod
 
class IDocumentProcessor(ABC):
    @abstractmethod
    def extract_text(self, file_path: str) -> str: ...
 
class DigitalPDFProcessor(IDocumentProcessor):
    def extract_text(self, file_path: str) -> str:
        # Uses PyMuPDF (fitz) to read native text layer
        ...
 
class ScannedPDFProcessor(IDocumentProcessor):
    def extract_text(self, file_path: str) -> str:
        # Converts pages to images, runs Tesseract OCR
        ...
 
class DocumentProcessorFactory:
    @staticmethod
    def get_processor(file_path: str) -> IDocumentProcessor:
        # Inspects file; returns DigitalPDFProcessor or ScannedPDFProcessor
        ...
\end{lstlisting}
 
% ── 3.2 ─────────────────────────────────────────────────────────────────────
\subsection{\texttt{normalization.py} --- Pre-Processing Layer (Decorator Pattern)}
 
\textbf{What it does.} Cleans the raw text returned by the extractor:
collapses whitespace, strips stray punctuation, and converts selected
word-form numbers (\emph{``twenty two''} $\to$ \texttt{22}) via a decorator.
 
\textbf{Design pattern: Decorator.} A base component returns text unchanged.
Each decorator wraps the previous layer, applies one specific transformation,
and passes the result forward. The application stacks them dynamically based
on which checkboxes the user enables in the UI.
 
\begin{mdframed}[style=patternbox]
\textbf{OOP Principle --- Single Responsibility (SRP).} No single class
performs more than one kind of text manipulation. Behaviours are composed
like Lego bricks rather than written as one monolithic \texttt{clean\_text()}
function.
\end{mdframed}
 
\begin{lstlisting}[caption={normalization.py --- interface and decorator skeleton}]
from abc import ABC, abstractmethod
 
class TextProcessor(ABC):
    @abstractmethod
    def process(self, text: str) -> str: ...
 
class BaseTextProcessor(TextProcessor):
    def process(self, text: str) -> str:
        return text   # foundation: pass text through unchanged
 
class WhitespaceRemover(TextProcessor):
    def __init__(self, wrapped: TextProcessor): self._wrapped = wrapped
    def process(self, text: str) -> str:
        # Collapse multiple spaces/newlines, then delegate upward
        ...
 
class PunctuationStripper(TextProcessor):
    def __init__(self, wrapped: TextProcessor): self._wrapped = wrapped
    def process(self, text: str) -> str:
        # Strip stray punctuation, then delegate upward
        ...
 
class WordToDigitConverter(TextProcessor):
    def __init__(self, wrapped: TextProcessor): self._wrapped = wrapped
    def process(self, text: str) -> str:
  # Convert "twenty two" -> "22" via word-map replacement, then delegate
        ...
\end{lstlisting}
 
A typical stack assembled in \texttt{app.py}:
\begin{lstlisting}[numbers=none]
pipeline = WordToDigitConverter(
               PunctuationStripper(
                   WhitespaceRemover(BaseTextProcessor())))
clean_text = pipeline.process(raw_text)
\end{lstlisting}
 
% ── 3.3 ─────────────────────────────────────────────────────────────────────
\subsection{\texttt{parser.py} --- Classification Layer (Polymorphism / LSP)}
 
\textbf{What it does.} Scans the cleaned text for numerical entities and
returns typed classification results across 24 labels (Date, Money, Percent,
Phone Number, License Plate, Time, Age, Temperature, Distance, Weight,
Height, Speed, Area, Volume, Percentage Change, Ratio, Range, Ordinal,
PIN Code, Account Number, Invoice ID, Tax ID, Quantity, Cardinal).
JSON serialization is now handled in a separate module.
 
\textbf{Design pattern: Polymorphism with Liskov Substitution.} An abstract
base class defines the contract. Three concrete engines implement it
differently --- Regex-only, spaCy NLP, and an LLM-backed parser.
The rest of the application calls \texttt{classify()} without knowing which
engine is active.
 
\begin{mdframed}[style=patternbox]
\textbf{OOP Principle --- Liskov Substitution (LSP) \& Encapsulation.}
\texttt{RegexParserClassifier}, \texttt{SpacyParserClassifier}, and
\texttt{LLMParserClassifier} are interchangeable. Switching engines requires
zero changes outside \texttt{parser.py} and UI configuration.
\end{mdframed}
 
\begin{lstlisting}[caption={parser.py --- abstract base and two concrete engines}]
from abc import ABC, abstractmethod
 
class BaseParserClassifier(ABC):
    @abstractmethod
  def classify(self, clean_text: str): ...
  # Returns a list of ClassificationResult dataclass objects
 
class RegexParserClassifier(BaseParserClassifier):
  def classify(self, clean_text: str):
    # Regex patterns for a 24-label numeric taxonomy
    ...

class SpacyParserClassifier(BaseParserClassifier):
  def classify(self, clean_text: str):
    # spaCy entities + matcher + regex fallback for full 24-label coverage
        ...
 
class LLMParserClassifier(BaseParserClassifier):
  def classify(self, clean_text: str):
    # Sends structured prompt with the 24-label set to LLM API;
    # maps response into ClassificationResult objects
        ...
\end{lstlisting}
 
\begin{table}[H]
\centering
\caption{Category mapping and normalisation examples}
\renewcommand{\arraystretch}{1.3}
\begin{tabular}{l l l}
\toprule
\textbf{Source} & \textbf{Category} & \textbf{original $\to$ normalised} \\
\midrule
spaCy \texttt{DATE}     & DATE           & ``14 August 2025'' $\to$ \texttt{14 August 2025} \\
spaCy \texttt{MONEY}    & MONEY          & ``Rs. 12,500'' $\to$ \texttt{Rs. 12,500} \\
spaCy \texttt{CARDINAL} & CARDINAL       & ``twenty'' $\to$ \texttt{twenty} \\
Regex match             & PHONE\_NUMBER & ``+91-98765-43210'' $\to$ \texttt{919876543210} \\
Regex match             & DATE           & ``14/08/2025'' $\to$ \texttt{14082025} \\
Regex match             & LICENSE\_PLATE & ``WB-24-X-9999'' $\to$ \texttt{WB24X9999} \\
\bottomrule
\end{tabular}
\end{table}
 
% ── 3.4 ─────────────────────────────────────────────────────────────────────
\subsection{\texttt{main.py} + \texttt{serialize.py} + \texttt{app.py} --- Orchestration, Serialization, and UI}

\textbf{What \texttt{main.py} does.} Defines \texttt{PipelineSettings},
\texttt{PipelineOutput}, and \texttt{ProcessingPipeline}. It performs
extractor selection, decorator stack construction, text normalization, and
classifier invocation.

\textbf{What \texttt{serialize.py} does.} Converts classifier output objects
into JSON-friendly records and strings, centralizing output formatting in one
place.
 
\textbf{What \texttt{app.py} does.} Provides the Streamlit UI: a file uploader, toggle
checkboxes for each normalisation decorator, and a sidebar to select the
classification engine (Regex, spaCy, or LLM). Displays normalized text and the
final JSON side-by-side; offers a \emph{Download JSON} button.
 
\textbf{OOP role: Thin Client + Service Layer.} \texttt{app.py} delegates core
pipeline behavior to \texttt{ProcessingPipeline} and output formatting to
\texttt{ResultSerializer}. This keeps UI and domain logic decoupled.
 
\begin{mdframed}[style=patternbox]
\textbf{OOP Principle --- Separation of Concerns.} The UI layer is completely
decoupled from extraction, normalisation, and classification logic. Any module
can be replaced with minimal impact on \texttt{app.py}.
\end{mdframed}
 
% ════════════════════════════════════════════════════════════════════════════
\section{End-to-End Data Flow}
% ════════════════════════════════════════════════════════════════════════════
 
\begin{figure}[H]
\centering
\begin{tikzpicture}[
  node distance=0.5cm and 0.35cm,
  every node/.style={font=\small},
  box/.style={rectangle, rounded corners=4pt, minimum width=2.5cm,
    minimum height=0.82cm, draw=AccentBlue, fill=LightGray,
    text=PrimaryBlue, align=center, thick},
  arr/.style={-Stealth, thick, color=AccentBlue},
  lbl/.style={font=\footnotesize\itshape\color{MidGray}},
]
\node[box] (pdf)  {PDF\\Input};
\node[box, right=of pdf]  (ext)  {extractor.py\\Factory};
\node[box, right=of ext]  (norm) {normalization.py\\Decorator Stack};
\node[box, right=of norm] (par)  {parser.py\\Classifier};
\node[box, right=of par]  (ser)  {serialize.py\\Serializer};
\node[box, right=of ser]  (out)  {JSON\\Output};
 
\draw[arr](pdf)--(ext);
\draw[arr](ext)--(norm)
\draw[arr](norm)--(par) 
\draw[arr](par)--(ser)
\draw[arr](ser)--(out) 
 
\node[lbl, below=0.2cm of ext]  {OCP};
\node[lbl, below=0.2cm of norm] {SRP};
\node[lbl, below=0.2cm of par]  {LSP};
\node[lbl, below=0.2cm of ser]  {SRP};
\end{tikzpicture}
\caption{End-to-end pipeline. Each module's governing OOP principle is shown below it.}
\end{figure}
 
\noindent The JSON output schema for every detected number is:
 
\begin{lstlisting}[language={}, numbers=none, caption={Required output schema (one entry per number)}]
{
  "original_text"    : "twenty-three August 2025",
  "normalized_value" : "2025-08-23",
  "category"         : "DATE"
}
\end{lstlisting}
 
% ════════════════════════════════════════════════════════════════════════════
\section{OOP Design Summary}
% ════════════════════════════════════════════════════════════════════════════
 
\begin{table}[H]
\centering
\renewcommand{\arraystretch}{1.35}
\begin{tabular}{l l l p{4.8cm}}
\toprule
\textbf{Module} & \textbf{Pattern} & \textbf{Principle} & \textbf{Benefit} \\
\midrule
\texttt{extractor.py}      & Factory Method & OCP & Add new file types without touching existing code \\
\texttt{normalization.py}  & Decorator      & SRP & Each transformation is isolated and composable \\
\texttt{parser.py}         & Polymorphism   & LSP & Regex, spaCy, and LLM engines are hot-swappable \\
\texttt{serialize.py}      & Utility / Adapter & SRP & Output schema formatting is centralized \\
\texttt{main.py}           & Service Layer  & Separation of Concerns & Pipeline orchestration is outside UI \\
\texttt{app.py}            & Client / Facade & Separation of Concerns & UI delegates domain logic to services \\
\bottomrule
\end{tabular}
\caption{Design patterns and OOP principles per module}
\end{table}

% ════════════════════════════════════════════════════════════════════════════
\section{Libraries Used}
% ════════════════════════════════════════════════════════════════════════════

\begin{table}[H]
\centering
\renewcommand{\arraystretch}{1.3}
\begin{tabular}{l p{5.2cm} p{4.8cm}}
\toprule
\textbf{Library} & \textbf{Purpose} & \textbf{Used In} \\
\midrule
\texttt{streamlit} & Web interface for file upload, options, and result display & \texttt{app.py} \\
\texttt{pymupdf (fitz)} & Native-text extraction from digital PDFs & \texttt{components/extractor.py} \\
\texttt{pdf2image} & Converts scanned PDF pages to images for OCR & \texttt{components/extractor.py} \\
\texttt{pytesseract} & OCR on scanned page images & \texttt{components/extractor.py} \\
\texttt{spacy} & NLP-based entity detection and fallback classification & \texttt{components/parser.py} \\
\texttt{openai} & LLM-backed classification for complex numeric contexts & \texttt{components/parser.py} \\
\texttt{num2words} & Number-to-word/word-to-number normalization helpers & \texttt{components/normalization.py} \\
\bottomrule
\end{tabular}
\caption{Primary external libraries used in the project pipeline}
\end{table}
 
% ════════════════════════════════════════════════════════════════════════════
\section{Conclusion}
% ════════════════════════════════════════════════════════════════════════════
 
The PDF Numerical Data Extractor satisfies the requirements of Question~5
while serving as a textbook demonstration of OOP design. The Factory Method
in \texttt{extractor.py} cleanly handles both digital and scanned PDFs; the
Decorator stack in \texttt{normalization.py} makes pre-processing steps
modular and user-configurable; Polymorphism in \texttt{parser.py} allows the
classification engine to be swapped between Regex, spaCy, and LLM with no
downstream changes. The \texttt{ProcessingPipeline} service in
\texttt{main.py} centralizes orchestration, and \texttt{serialize.py}
encapsulates JSON output construction. The Streamlit front-end in
\texttt{app.py} remains a thin client layer. Together, the modules
demonstrate OCP, SRP, LSP, Encapsulation, and Separation of Concerns across a
single, cohesive pipeline.
 
% ─── Bibliography ───────────────────────────────────────────────────────────
\begin{thebibliography}{9}
\bibitem{spacy}
Honnibal et al.\ (2020). \textit{spaCy: Industrial-strength NLP.}
\url{https://spacy.io}
\bibitem{tesseract}
Smith, R.\ (2007). An overview of the Tesseract OCR engine.
\textit{ICDAR 2007.} IEEE.
\bibitem{pymupdf}
Artifex Software. \textit{PyMuPDF Documentation.}
\url{https://pymupdf.readthedocs.io}
\bibitem{gamma}
Gamma et al.\ (1994). \textit{Design Patterns: Elements of Reusable
Object-Oriented Software.} Addison-Wesley.
\end{thebibliography}
 
\end{document}